
\begin{theorem}[Sectorial Hamilton decomposition of the geometric branch]
\label{thm:sectorial-hamilton-decomposition}
On the admissible Lorentzian reconstruction, the full Hamiltonian splits as
\[
H_{\mathrm{full}}
:=
H_{\mathrm{matt}}\widehat\otimes 1
+
1\widehat\otimes H_{\mathrm{grav}}
+
H_{\mathrm{int}},
\qquad
H_{\mathrm{grav}}=H_{2}\oplus H_{\mathrm{sc}}.
\]
The gap / clustering closure applies to the matter block and the massive scalaron block. The
transverse-traceless graviton block is massless, but it is controlled by Osterwalder--Schrader
reflection positivity together with the positive TT quadratic form of
\Cref{thm:gravity-cosmology-closure}.
\end{theorem}

\begin{proof}[Proof sketch]
The geometric Hodge theorem already identifies the physical Hilbert-space factorization
$\mathcal H_{\mathrm{phys}}\cong\mathcal H_2\otimes\mathcal H_{\mathrm{sc}}\otimes\mathcal H_{\mathrm{matt}}$.
Writing the reconstructed Hamiltonian relative to that factorization isolates the massless
transverse-traceless graviton sector from the massive scalaron / matter sectors. Hence the gap
statements later used for admissible vacuum closure are not assigned to the graviton block; that
block is controlled instead by reflection positivity and the TT positivity statement already proved
on the geometric branch.
\end{proof}

\begin{theorem}[Unique admissible UV profile fixed point]
\label{thm:geometric-uv-fixed-point}
Let
\[
\mathcal M_{\mathrm{adm}}^{\mathrm{br}}
:=
\left\{
\nu\ge 0\ \text{finite on }[0,\infty)\;:\;
f_\nu(t):=\int_0^\infty e^{-st}\,d\nu(s)\ \text{is branch compatible}
\right\}
\]
be the admissible Bernstein cone written on the measure side. Fix the branch-matching slice by the
closed low-curvature moments reproducing the Einstein, Higgs, and $R^2$ coefficients. Define the
normalized positive RG map
\[
(T\nu)(dt)
:=
\frac{\int_0^\infty K(s,t)\,d\nu(s)}{\int_0^\infty \psi(s)\,d\nu(s)}\,dt,
\qquad
\psi(s)>0.
\]
Let $\mathcal R$ denote the induced map on Laplace profiles,
\[
\mathcal R[f_\nu]:=f_{T\nu}.
\]
Assume that on the branch-matching slice the kernel satisfies the uniform bounds
\[
0<m_{\mathrm{ker}}\le K(s,t)\le M_{\mathrm{ker}}<\infty.
\]
Then $T$ has finite Hilbert-projective diameter and obeys the Birkhoff contraction estimate
\[
d_H(T\nu_1,T\nu_2)
\le
\eta\,d_H(\nu_1,\nu_2),
\qquad
\eta\le
\tanh\!\left(\frac14\log\frac{M_{\mathrm{ker}}}{m_{\mathrm{ker}}}\right)<1.
\]
Consequently there exists a unique admissible fixed measure $\nu_\star$ and therefore a unique
admissible fixed profile
\[
\mathcal R[f_\star]=f_\star,
\qquad
f_\star(t)=\int_0^\infty e^{-st}\,d\nu_\star(s),
\]
whose first moments reproduce the Einstein, Higgs, and $R^2$ coefficients of the closed branch.
\end{theorem}

\begin{proof}
The positive measure cone is the natural projective cone for admissible Bernstein profiles, and the
branch-matching moment conditions cut out a normalized slice preserving the closed low-curvature
coefficients
\[
(\Mpl^2,Z_H,M^2),
\]
so any admissible profile on that slice has the same first moments as the closed branch. Because
$K$ and $\psi$ are positive, $T$ preserves positivity and the matching slice. The uniform kernel
bounds imply a finite Hilbert-projective diameter for the image of that slice; explicitly, the
projective diameter is bounded by $\log(M_{\mathrm{ker}}/m_{\mathrm{ker}})$. Birkhoff's theorem for
positive kernel operators therefore gives the displayed contraction factor
\[
\eta\le \tanh\!\left(\frac14\log\frac{M_{\mathrm{ker}}}{m_{\mathrm{ker}}}\right)<1.
\]
Hence $T$ has a unique fixed projective ray on the normalized slice. The normalization by
$\int\psi\,d\nu$ fixes the scale, so the fixed ray is represented by one actual measure
\[
\nu_\star\in\mathcal M_{\mathrm{adm}}^{\mathrm{br}}.
\]
Its Laplace profile
\[
f_\star(t)=\int_0^\infty e^{-st}\,d\nu_\star(s)
\]
is the unique admissible UV profile compatible with the closed branch, and its first moments are
exactly the prescribed Einstein, Higgs, and $R^2$ coefficients.
\end{proof}

\begin{theorem}[Constructive geometric measure on the admissible branch]
\label{thm:constructive-geometric-measure}
Let $(V_n)_{n\ge 1}$ be a blocked geometric exhaustion and define the finite-volume geometric
measures
\[
:=
d\mu_n
\propto
\exp\!\bigl(-\tr f_\star(D_{V_n}^2/\Lambda^2)\bigr)\,d\nu_{V_n}
\]
on the BRST-reduced admissible geometric configuration space. Then the family $(d\mu_n)$ is
projectively consistent, reflection positive, and converges along the exhaustion to a projective
limit measure. In particular the geometric partition function is defined constructively by
\[
Z_{\mathrm{grav}}
=
\lim_{n\to\infty}
\int e^{-\tr f_\star(D_{V_n}^2/\Lambda^2)}\,d\nu_{V_n}.
\]
\end{theorem}

\begin{proof}
For every blocked volume $V_n$, the weight
\[
\exp\!\bigl(-\tr f_\star(D_{V_n}^2/\Lambda^2)\bigr)
\]
is well defined because $f_\star$ is positive on the admissible Bernstein cone. If
\[
V_n\subset V_m,
\]
restriction of geometric configurations from $V_m$ to $V_n$ respects the BRST-reduced admissible
action because the blocked Dirac operators and the fixed profile $f_\star$ were constructed to be
compatible under exhaustion. This is exactly the projective-consistency statement.

Reflection positivity is inherited from the admissible/geometric OS structure established earlier:
the reflected quadratic forms remain nonnegative on every blocked volume, and the same property is
preserved by the positive weight associated with $f_\star$. Hence each $d\mu_n$ is reflection
positive.

Finally, the admissible positivity and coercivity bounds prevent escape of mass along the blocked
exhaustion, so the projectively consistent family $(d\mu_n)$ admits a projective-limit measure. The
geometric partition function is therefore defined constructively by the stated limit, rather than
only heuristically read off from low-curvature coefficients.
\end{proof}

\begin{theorem}[Spectral UV completion on the admissible branch]
\label{thm:spectral-uv-completion}
Let $f_\star$ be the unique admissible profile of \Cref{thm:geometric-uv-fixed-point}, and define
the transverse-traceless and scalaron form factors in Stieltjes form
\[
\mathcal F_\star(z)
:=
a_F+\int_0^\infty \frac{d\mu_F(t)}{t+z},
\qquad
\mathcal G_\star(z)
:=
a_G+\int_0^\infty \frac{d\mu_G(t)}{t+z},
\]
with
\[
a_F,a_G\ge 0,
\qquad
\mu_F,\mu_G\ge 0.
\]
Then the Euclidean two-point gravitational kernel on the admissible branch takes the spectral form
\[
\Gamma_{\mathrm{grav}}^{(2)}(p^2)
=
p^2\,\mathcal F_\star\!\left(\frac{p^2}{\Lambda_{\mathrm{UV}}^2}\right)\Pi_{\mathrm{TT}}
+
\bigl(p^2+M^2\bigr)\mathcal G_\star\!\left(\frac{p^2}{\Lambda_{\mathrm{UV}}^2}\right)\Pi_{\mathrm{sc}}.
\]
\begin{enumerate}[label=(\roman*), leftmargin=1.5em, topsep=3pt, itemsep=2pt]
    \item $\mathcal F_\star$ and $\mathcal G_\star$ are holomorphic and zero-free on
    $\mathbb C\setminus(-\infty,0]$;
    \item the only physical poles are the massless graviton pole and the positive scalaron pole;
    \item no ghost pole occurs on the admissible branch;
    \item the constructive geometric measure of \Cref{thm:constructive-geometric-measure} supplies the
    corresponding UV correlators as moments of one and the same admissible measure.
\end{enumerate}
Thus the local Einstein plus $R^2$ action is the low-curvature expansion of a spectral UV object
rather than the full ultraviolet definition of the branch.
\end{theorem}

\begin{proof}
The second variation of the fixed-profile spectral action on the BRST-reduced geometric sector
produces positive spectral measures on the transverse-traceless and scalar blocks; these are
recorded above as the Stieltjes data $(a_F,\mu_F)$ and $(a_G,\mu_G)$. The projector decomposition
\[
\Pi_{\mathrm{TT}}+\Pi_{\mathrm{sc}}
\]
separates the massless graviton and scalaron channels, yielding the displayed two-point kernel.

If $\Im z>0$, then for every $t\ge 0$,
\[
\Im\frac{1}{t+z}
=
-\frac{\Im z}{(t+\Re z)^2+(\Im z)^2}
<0.
\]
Hence
\[
\Im \mathcal F_\star(z)<0,
\qquad
\Im \mathcal G_\star(z)<0
\qquad
(\Im z>0),
\]
so neither form factor can vanish in the upper half-plane. By conjugation symmetry the same holds
in the lower half-plane. On the positive real axis the Stieltjes integrands are strictly positive,
and the branch-matching coefficients force the resulting values to remain positive. Therefore
$\mathcal F_\star$ and $\mathcal G_\star$ are zero-free on
$\mathbb C\setminus(-\infty,0]$.

The inverse propagators can therefore acquire no additional pole beyond the physical roots at
\[
p^2=0
\qquad\text{and}\qquad
p^2=-M^2.
\]
Hence there is no ghost pole on the admissible branch. Finally,
\Cref{thm:constructive-geometric-measure} provides one projective-limit measure with fixed profile
$f_\star$; its moments are precisely the UV correlators of this same spectral kernel. Therefore the
Einstein plus $R^2$ action obtained earlier is only the first low-curvature expansion of the
spectral UV completion.
\end{proof}

The E$_8$ scale ladder is not used as a primitive input in the main text. Its arithmetic
continuations are deferred to Appendix~\ref{app:constants}.

\section{Closure architecture B: Yukawas, masses, flavor, hadronic admissibility, and the strong-CP sector}

The second closure-architecture block is the transport closure. Its role is to turn the carrier
packet and the family/occupancy closure into actual masses, mixings, hadronic admissibility, and
control of the strong-CP sector. The logic is again staged: first the discrete phase operator and
the transport kernel, then residue and winding data, then matrix-factorized Yukawas, hadronic
selection, determinant-line phase extraction, and finally the strong-CP closure theorem.

\subsection{Discrete phase operator and transport grammar}

Let
\[
\omega=e^{2\pi i/3},
\qquad
U_6=\mathrm{diag}(1,\omega,\omega^2,-1,-\omega,-\omega^2).
\]
This makes the phase grammar concrete rather than rhetorical. The six admissible phases are shown as a finite hexagonal transport geometry in \Cref{fig:z6-circle}: the cusp values live on the same real axis as the positive-real phase, the admissible gap is visible directly, and the holonomy phase can be read as a triangle on the same discrete transport base.

\begin{remark}[Hexagonal spectral geometry]
The operator $U_6$ is a unitary circulant with eigenphases $e^{in\pi/3}$ for $n=0,\dots,5$ and is unitarily equivalent to the one-step translation operator on a six-site ring. The transport kernel $D_y=y\,\mathbf 1-\delta_{\mathrm{ph}}U_6$ is therefore the resolvent denominator of a finite quantum walk, equivalently a tight-binding transfer problem, on the discrete $C_6$ cycle. The regularized inverse introduced below plays the role of the corresponding Green kernel. CKM and PMNS mixings, mass hierarchies, and the Möbius quotients of \Cref{eq:mobius-quotient} are thus not loose analytic fits but spectral-theory statements on a regular hexagon. The near-pole hierarchy at $y\approx\delta_{\mathrm{ph}}$ becomes the resonance of a hexagonal resolvent rather than tuned parameter choice.
\end{remark}

\begin{figure}[t]
\centering
\scriptsize
\begin{tikzpicture}[scale=1.1,>=Latex]
\coordinate (O) at (0,0);
\foreach \name/\ang in {p0/0,p1/60,p2/120,p3/180,p4/240,p5/300} {
    \coordinate (\name) at ({2*cos(\ang)},{2*sin(\ang)});
}
\draw[thick, black!70] (p0) -- (p1) -- (p2) -- (p3) -- (p4) -- (p5) -- cycle;
\draw[gray!25, dashed] (O) -- (p0);
\draw[gray!25, dashed] (O) -- (p1);
\draw[gray!25, dashed] (O) -- (p2);
\draw[gray!25, dashed] (O) -- (p3);
\draw[gray!25, dashed] (O) -- (p4);
\draw[gray!25, dashed] (O) -- (p5);
\foreach \P in {p0,p1,p2,p3,p4,p5} {
    \filldraw[blue!60!black] (\P) circle (2pt);
}
\node[anchor=west, font=\scriptsize] at (2.18,0.12) {$1$};
\node[anchor=south west, font=\scriptsize] at (1.1,1.85) {$e^{i\pi/3}$};
\node[anchor=south east, font=\scriptsize] at (-1.1,1.85) {$e^{2i\pi/3}$};
\node[anchor=east, font=\scriptsize] at (-2.18,0.12) {$-1$};
\node[anchor=north east, font=\scriptsize] at (-1.1,-1.85) {$e^{-2i\pi/3}$};
\node[anchor=north west, font=\scriptsize] at (1.1,-1.85) {$e^{-i\pi/3}$};

\draw[->, thick, red!65!black] (1.15,1.05) arc[start angle=42,end angle=-42,radius=1.2];
\node[red!65!black, font=\scriptsize, anchor=west] at (2.55,0.25) {CP: $\theta\mapsto-\theta$};

\filldraw[orange!80!black] (0.67,0) circle (2pt);
\filldraw[orange!80!black] (1.33,0) circle (2pt);
\filldraw[green!60!black] (1.1998,0) circle (2.4pt);
\node[below, font=\scriptsize] at (0.67,-0.08) {$\frac13$};
\node[below, font=\scriptsize] at (1.33,-0.08) {$\frac23$};
\node[above, font=\scriptsize] at (1.1998,0.08) {$\delta_{\mathrm{ph}}$};
\draw[<->, thick, black!65] (1.1998,-0.32) -- (1.33,-0.32);
\node[below, font=\tiny] at (1.265,-0.34) {$\epsilon_f$};
\node[align=left, anchor=west, font=\scriptsize, text=black!70] at (2.55,-0.7)
    {positive-real branch\\near-pole lock at $y\approx\delta_{\mathrm{ph}}$\\gap $\epsilon_f=\frac23-\delta_{\mathrm{ph}}$};
\end{tikzpicture}
\caption{Finite transport geometry on the hexagon $C_6$. The six $U_6$ phases form the discrete
transport base, while the admissible cusps $y\in\{1,\frac13,\frac23\}$ and the branch pole
$\delta_{\mathrm{ph}}$ lie on the positive-real axis with gap
$\epsilon_f=\frac23-\delta_{\mathrm{ph}}$.}
\label{fig:z6-circle}
\end{figure}

The family-side holonomy phase is separated in \Cref{fig:holonomy-triangle} so that the transport
base and the flavor readout remain visually distinct.

\begin{figure}[t]
\centering
\scriptsize
\begin{tikzpicture}[>=Latex]
    \node[draw=black!65, fill=gray!10, rounded corners, text width=4.8cm, minimum height=0.85cm, align=center, font=\bfseries] at (0,1.15)
        {family-side holonomy on the transport grammar};
    \coordinate (A) at (0,0.1);
    \coordinate (B) at (-1.45,-1.8);
    \coordinate (C) at (1.45,-1.8);
    \draw[thick, orange!75!black] (A) -- (B) -- (C) -- cycle;
    \filldraw[orange!75!black] (A) circle (2pt) (B) circle (2pt) (C) circle (2pt);
    \node[above, font=\scriptsize] at (A) {$1$};
    \node[left, font=\scriptsize] at (B) {$2$};
    \node[right, font=\scriptsize] at (C) {$3$};
    \draw[->, thick, blue!60!black] (-0.2,-0.45) arc[start angle=110,end angle=15,radius=0.95];
    \node[draw=green!50!black, fill=green!8, rounded corners, text width=3.7cm, minimum height=0.95cm, align=center] at (0,-3.15)
        {$\delta_{\mathrm{CKM}}=\arg\tr\Pi_\triangle(\nabla_F^\star)$\\oriented holonomy triangle};
\end{tikzpicture}
\caption{Holonomy triangle on the same discrete transport grammar. The CKM phase is read as the
oriented three-cycle phase $\delta_{\mathrm{CKM}}=\arg\tr\Pi_\triangle(\nabla_F^\star)$ rather
than as an additional transport parameter.}
\label{fig:holonomy-triangle}
\end{figure}

